\hnheader[Rückwärts durch das Netz: die Kettenregel, systematisch angewandt.][Der Gradient kostet nur so viel wie der Forward-Pass: $O(N)$!]{Backpropagation}{\& Optimierung}{Gradienten effizient berechnen und nutzen}
\hnsrc{main\_notes.pdf (Backpropagation), notes/cs229-notes-backprop.pdf}

\begin{hngrid}
%
\begin{hnp}[raster multicolumn=2,acc=hnorange]{1}{Theorem}
Lässt sich der Verlust in einem Forward-Pass mit Aufwand $N$ berechnen, so lässt sich der \emph{vollständige} Gradient $\nabla_\theta J$ ebenfalls in $O(N)$ berechnen. Backprop ist ``nur'' die Kettenregel mit Zwischenspeicherung.
\end{hnp}
%
\begin{hnp}[raster multicolumn=2,acc=hnpurple]{2}{Formeln}
\hnleg{$u^{[\ell]}$ Zwischenwert nach Modul $\ell$, $M_\ell$ Modul (Schicht), $\theta^{[\ell]}$ dessen Parameter, $J$ Verlust, $B$ Rückwärtsfunktion (transponierte Ableitung des Moduls), $L$ Anzahl Module.}
Forward: $u^{[\ell]}=M_\ell(u^{[\ell-1]};\theta^{[\ell]})$, $J=u^{[L]}$.\\
Backward:
\begin{align*}
\tfrac{\partial J}{\partial u^{[L]}}&=1\\
\tfrac{\partial J}{\partial u^{[\ell-1]}}&=B\bigl[M_\ell,u^{[\ell-1]}\bigr]\bigl(\tfrac{\partial J}{\partial u^{[\ell]}}\bigr)\\
\tfrac{\partial J}{\partial\theta^{[\ell]}}&=B_\theta\bigl[M_\ell,\theta^{[\ell]}\bigr]\bigl(\tfrac{\partial J}{\partial u^{[\ell]}}\bigr)
\end{align*}
\end{hnp}
%
\begin{hnp}[raster multicolumn=2,acc=hnteal]{3}{Skizze: Vor- \& Rückwärts}
\centering
\begin{tikzpicture}[scale=1.35,>=Stealth,every node/.style={draw=hnnavy,fill=hnblue!50,rounded corners=1pt,font=\tiny,inner sep=2.5pt}]
  \node (x) at (0,0) {$x$};
  \node (a) at (1.2,0) {$M_1$};
  \node (b) at (2.4,0) {$M_2$};
  \node[fill=hnpink] (j) at (3.6,0) {$J$};
  \draw[->,hnnavy,thick] (x)--(a); \draw[->,hnnavy,thick] (a)--(b); \draw[->,hnnavy,thick] (b)--(j);
  \draw[->,hnorange,thick] (j.south) to[bend left=25] node[draw=none,fill=none,below,font=\tiny]{$\partial J/\partial u^{[2]}$} (b.south);
  \draw[->,hnorange,thick] (b.south) to[bend left=25] (a.south);
  \node[draw=none,fill=none,hnnavy] at (1.8,.55) {Forward};
  \node[draw=none,fill=none,hnorange] at (1.8,-1.0) {Backward};
\end{tikzpicture}
\end{hnp}
%
\begin{hnp}[raster multicolumn=3,acc=hnnavy]{4}{Backprop im MLP}
\begin{pcode}[Backpropagation]
\cm{\# Forward (Zwischenwerte speichern)}\\
$a_0\leftarrow x$;\; \kw{for} $\ell=1..L$:\\
\hii $z_\ell\leftarrow W_\ell a_{\ell-1}+b_\ell$;\; $a_\ell\leftarrow\sigma(z_\ell)$\\
\cm{\# Backward}\\
$\delta_L\leftarrow\partial J/\partial z_L$\\
\kw{for} $\ell=L..1$:\\
\hii $\partial J/\partial W_\ell\leftarrow\delta_\ell\,a_{\ell-1}^{\top}$;\; $\partial J/\partial b_\ell\leftarrow\delta_\ell$\\
\hii $\delta_{\ell-1}\leftarrow(W_\ell^{\top}\delta_\ell)\odot\sigma'(z_{\ell-1})$
\end{pcode}
\end{hnp}
%
\begin{hnp}[raster multicolumn=3,acc=hngreen]{5}{Mini-Batch-SGD}
\begin{pcode}[Mini-Batch SGD]
\kw{init} $\theta$ (zufällig)\\
\kw{repeat} bis Konvergenz:\\
\hii Batch $\mathcal B\subset\{1..n\}$, $|\mathcal B|=B$\\
\hii $g\leftarrow\frac1B\sum_{i\in\mathcal B}\nabla_\theta J^{(i)}(\theta)$\\
\hii $\theta\leftarrow\theta-\alpha\,g$
\end{pcode}
($\mathcal B$ Batch, $B$ Batchgröße, $J^{(i)}$ Verlust von Beispiel $i$, $g$ Gradientenschätzung, $\alpha$ Lernrate.) $B=1$: reines SGD (verrauscht); $B=n$: Batch-GD (teuer). Kleine Batches sind billig \emph{und} wirken leicht regularisierend.\\
\hlt{Web-Zusatz:} in der Praxis oft Momentum bzw.\ Adam statt reinem SGD.
\end{hnp}
%
\begin{hnex}[raster multicolumn=3]{Beispiel: Backprop von Hand}
$\hat y=w_2\,\mathrm{ReLU}(w_1x)$, $x=2$, $y=0$, $w_1=1$, $w_2=0.5$, $J=\frac12(\hat y-y)^2$.\\
\hnsol\ Forward: $z=2$, $a=2$, $\hat y=1$, $J=0.5$.\\
Backward: $\partial J/\partial\hat y=1$; $\partial J/\partial w_2=1\cdot a=2$; $\partial J/\partial a=1\cdot w_2=0.5$; $\partial J/\partial z=0.5$ (ReLU$'{=}1$); $\partial J/\partial w_1=0.5\cdot x=1$.\\
SGD ($\alpha{=}0.1$): $w_2=0.3$, $w_1=0.9\Rightarrow\hat y=0.54$, $J=\ora{0.146}<0.5$ \hlm{Verlust sinkt}.
\end{hnex}
%
\begin{hnp}[raster multicolumn=3,acc=hnrose]{6}{Praxis \& Fallstricke}
\begin{itemize}
\item \textbf{Gradient-Check}: $\frac{J(\theta+\varepsilon)-J(\theta-\varepsilon)}{2\varepsilon}$ mit Backprop vergleichen.
\item \textbf{Vanishing/Exploding Gradients}: Sigmoid/Tanh sättigen; Abhilfe: ReLU, ResNet, LayerNorm, gute Initialisierung.
\item Lernrate zu groß $\to$ Divergenz, zu klein $\to$ Stillstand.
\end{itemize}
\end{hnp}
%
\begin{hnp}[raster multicolumn=2,acc=hngreen]{7}{Vorteile}
\begin{hnpro}
\item exakter Gradient
\item Aufwand $\approx$ ein Forward-Pass
\item Autodiff-Frameworks übernehmen es
\end{hnpro}
\end{hnp}
%
\begin{hnp}[raster multicolumn=2,acc=hnred]{8}{Grenzen}
\begin{hncon}
\item Zwischenwerte brauchen Speicher
\item nichtkonvex: nur lokale Minima
\item Hyperparameter-Tuning nötig
\end{hncon}
\end{hnp}
%
\begin{hnp}[raster multicolumn=2,acc=hnorange]{9}{Key Takeaways}
\begin{hnchk}
\item Backprop = Kettenregel rückwärts
\item Kosten $O(N)$
\item SGD in Mini-Batches
\end{hnchk}
\end{hnp}
%
\begin{hnp}[raster multicolumn=6,acc=hnpurple,colback=hnlilac!30]{?}{Selbsttest: kannst du das erklären?}
\hnquiz{Aufwand von Backprop?}{$O(N)$, wie ein Forward-Pass.}{Was ist ein Gradient-Check?}{$\frac{J(\theta+\varepsilon)-J(\theta-\varepsilon)}{2\varepsilon}$ mit Backprop vergleichen.}{Wirkung der Batchgröße?}{$B{=}1$ verrauscht, $B{=}n$ teuer; Mini-Batch als Kompromiss.}
\end{hnp}
%
\begin{hnbanner}[raster multicolumn=6]
„Rückwärts rechnen, vorwärts lernen.“
\end{hnbanner}
\end{hngrid}
